\chapter{How to Perform CPR}

\section*{Definition and Overview}
Cardiopulmonary resuscitation (CPR) is a life-saving first aid technique used when someone's heart has stopped or they are not breathing. CPR combines chest compressions, which manually pump blood to the brain and vital organs, with rescue breaths that provide oxygen. Starting CPR immediately can double or triple a person's chance of survival after cardiac arrest. The current Australian recommendation is compressions at a rate of 100--120 per minute, to a depth of about 5 cm in adults, with minimal interruption. Everyone can learn CPR, and doing something is always better than doing nothing.

\section*{Epidemiology}
Around 20,000--30,000 Australians suffer an out-of-hospital cardiac arrest each year, and survival is only around 10\%. Bystander CPR is the single most important factor in survival: it is performed in only about half of witnessed arrests, and increasing bystander CPR rates would save hundreds of lives annually. Survival falls by about 10\% for every minute without CPR or defibrillation. Community CPR training and publicly accessible defibrillators (AEDs) have improved outcomes in many areas.

\section*{Aetiology and Pathophysiology}
\begin{itemize}
    \item \textbf{Cardiac arrest:} the heart stops pumping effectively, usually from a heart attack, arrhythmia, drowning, drug overdose, or trauma
    \item \textbf{Without CPR:} the brain is damaged within minutes as oxygen is depleted
    \item \textbf{Chest compressions:} manually compress the heart between the breastbone and spine, pushing blood to the brain and organs
    \item \textbf{Rescue breaths:} provide oxygen to maintain organ function
    \item \textbf{Defibrillation:} an electrical shock can restore a normal rhythm in shockable arrests
    \item \textbf{Chain of survival:} early recognition and call for help, early CPR, early defibrillation, and early advanced care
    \item \textbf{Pathophysiology:} effective CPR maintains some circulation until the heart can be restarted
\end{itemize}

\section*{Clinical Features}
\begin{itemize}
    \item The person is unresponsive and not breathing normally (or only gasping)
    \item No response to calling or gentle shaking
    \item No normal breathing within 10 seconds of assessment
    \item CPR is started immediately and continued until help arrives
    \item An AED is used as soon as one is available
    \item CPR continues with minimal interruption
\end{itemize}

\section*{Differential Diagnosis}
\begin{itemize}
    \item Unresponsiveness from fainting, seizure, or stroke may include abnormal breathing but is not cardiac arrest; CPR is only for no normal breathing
    \item Agonal gasps are not normal breathing and indicate cardiac arrest
    \item Opioid overdose causes slow, shallow breathing that may stop; naloxone and ventilation are key
    \item When in doubt, start CPR: the risk of harm is far lower than the risk of inaction
\end{itemize}

\section*{When to Seek Medical Care}
Call 000 immediately for any unresponsive person who is not breathing normally, put the phone on speaker, and follow the dispatcher's instructions.

\textbf{Call 000 (or local emergency services) for:}
\begin{itemize}
    \item Any suspected cardiac arrest
    \item Unresponsiveness with abnormal or absent breathing
    \item Drowning, choking, drug overdose, or severe trauma with collapse
    \item A person who stops responding after an injury or illness
    \item Continue CPR until emergency services arrive
\end{itemize}

\section*{Diagnosis and Investigations}
\begin{itemize}
    \item \textbf{Recognition:} check for response and normal breathing for up to 10 seconds
    \item \textbf{Call for help:} send someone to call 000 and fetch an AED
    \item \textbf{CPR:} 30 compressions then 2 breaths, at 100--120 compressions per minute
    \item \textbf{AED use:} turn it on and follow the voice prompts; shocks restore rhythm in shockable arrests
    \item \textbf{Minimal interruption:} limit pauses in compressions
    \item \textbf{Aftercare:} hand over to ambulance crews and follow their instructions
\end{itemize}

\section*{Management}
\begin{itemize}
    \item \textbf{Adult CPR:} place hands on the centre of the chest, push hard and fast (about 5 cm deep, 100--120 per minute), allow the chest to recoil
    \item \textbf{Compression-only CPR:} for adults who are unwilling or unable to give breaths, continuous compressions are effective
    \item \textbf{Child and infant CPR:} compress about one-third of the chest depth, and give gentler rescue breaths
    \item \textbf{Drowning and children:} begin with rescue breaths because oxygen deprivation is the cause
    \item \textbf{AED:} attach pads, follow prompts, and resume CPR after any shock
    \item \textbf{Recovery position:} for an unresponsive person who is breathing normally
    \item \textbf{Training:} formal CPR courses build skill and confidence
\end{itemize}

\section*{Complications}
\begin{itemize}
    \item Delayed CPR reduces survival substantially
    \item Rib fractures from compressions are common but acceptable given the alternative
    \item Ineffective compressions (too slow or shallow) reduce blood flow
    \item Prolonged interruption of compressions worsens outcomes
    \item Failing to use a nearby AED when available
    \item Hesitation from fear of doing harm, when action would have helped
\end{itemize}

\section*{Prognosis and Outcomes}
Bystander CPR and early defibrillation dramatically improve survival from cardiac arrest, with survival rates up to three times higher when CPR is started immediately and an AED is used early. Most survivors have good neurological recovery when CPR is prompt and effective. Even imperfect CPR is far better than none. With training and willingness to act, community members are the first and most important link in the chain of survival.

\section*{Prevention and Self-Care}
\begin{itemize}
    \item Learn CPR through a certified course and refresh skills regularly
    \item Learn to use an AED
    \item Know the location of AEDs in your workplace and community
    \item Recognise the warning signs of heart attack and call 000 early
    \item Encourage family and colleagues to learn CPR
    \item Practise the DRSABCD action plan: Danger, Response, Send for help, Airway, Breathing, CPR, Defibrillation
    \item Stay updated with Australian Resuscitation Council guidelines
    \item Seek help for anxiety about performing first aid by building skills and confidence
\end{itemize}

\section*{Sources and Further Reading}
The following reputable sources provide further information for healthcare students and patients seeking reliable, current advice about CPR.

\begin{itemize}
    \item Australian Resuscitation Council. \textit{CPR guidelines and first aid.} https://resus.org.au/
    \item St John Ambulance Australia. \textit{CPR courses and first aid information.} https://stjohn.org.au/
    \item Heart Foundation. \textit{CPR and defibrillation information.} https://www.heartfoundation.org.au/
    \item Healthdirect. \textit{How to perform CPR.} https://www.healthdirect.gov.au/cpr
    \item Red Cross Australia. \textit{First aid and CPR training.} https://www.redcross.org.au/
    \item British Heart Foundation. \textit{Hands-only CPR guide.} https://www.bhf.org.uk/
\end{itemize}
